% =============================================================================
% Module 15: Solutions to Exercises
% =============================================================================
% This file is conditionally included from main.tex via \ifsolutions
% Every answer is attested by Solutions/15_PowerLaw_PointMass/problems_15.wl
% =============================================================================

\section{Solutions to Exercises}
\label{sec:pl_solutions}

\noindent
Full worked solutions are also available in the Mathematica script
\solnlink{15_PowerLaw_PointMass/problems\_15.wl}{problems\_15.wl},
which verifies each answer symbolically and numerically.

% ----- Exercise: kappa by inversion -----
\subsection*{Exercise~\ref{ex:pl_kappa_inversion}: Power-Law Convergence by Inversion}

Form the combination $\theta\,\alpha_{\rm PL}=\thetaE^{\gp-1}\theta^{3-\gp}$ and differentiate:
\begin{equation*}
    \deriv{}{\theta}\big[\theta\,\alpha_{\rm PL}\big]=(3-\gp)\,\thetaE^{\gp-1}\theta^{2-\gp} .
\end{equation*}
Dividing by $2\theta$,
\begin{equation*}
    \kappab(\theta)=\frac{1}{2\theta}\deriv{}{\theta}\big[\theta\,\alpha_{\rm PL}\big]
      =\frac{3-\gp}{2}\,\thetaE^{\gp-1}\theta^{1-\gp}
      =\frac{3-\gp}{2}\left(\frac{\thetaE}{\theta}\right)^{\gp-1} .
\end{equation*}
In the other direction,
\begin{equation*}
    \frac{2}{\theta}\int_0^\theta\kappab(\theta')\,\theta'\,\dd\theta'
    =\frac{2}{\theta}\cdot\frac{3-\gp}{2}\thetaE^{\gp-1}\int_0^\theta\theta'^{\,2-\gp}\dd\theta'
    =\frac{2}{\theta}\cdot\frac{3-\gp}{2}\thetaE^{\gp-1}\frac{\theta^{3-\gp}}{3-\gp}
    =\thetaE^{\gp-1}\theta^{2-\gp} ,
\end{equation*}
where the integral converges at the origin only for $\gp<3$ (finite
projected mass).  Finally $\bar\kappa(\theta)=\alpha/\theta=(\thetaE/\theta)^{\gp-1}$,
so $\bar\kappa(\thetaE)=1$ for every slope: the Einstein radius is
where the mean interior convergence is unity.
\emph{[Attested: \solnlink{15_PowerLaw_PointMass/problems\_15.wl}{problems\_15.wl}, checks 1a--1d.]}

% ----- Exercise: point-mass magnification -----
\subsection*{Exercise~\ref{ex:pl_point_mass_mag}: Point-Mass Magnification}

With $\alpha_\bullet=\thetaBH^2/\theta$ we have
$\lambda_t=1-\thetaBH^2/\theta^2$ and $\lambda_r=1+\thetaBH^2/\theta^2$, so
\begin{equation*}
    \mu=\frac{1}{\lambda_t\lambda_r}
      =\frac{1}{(1-\thetaBH^2/\theta^2)(1+\thetaBH^2/\theta^2)}
      =\frac{1}{1-\thetaBH^4/\theta^4}
      =\frac{\theta^4}{\theta^4-\thetaBH^4} .
\end{equation*}
Equivalently, $\mu=(\theta/\ysrc)(\dd\theta/\dd\ysrc)$ for the map
$\ysrc=\theta-\thetaBH^2/\theta$.  Since $\lambda_r>0$ always, the
parity is set by $\lambda_t$: positive outside $\thetaBH$, negative
inside.

The lens equation is $\theta^2-\ysrc\,\theta-\thetaBH^2=0$, with roots
$\theta_\pm=\tfrac12(\ysrc\pm\sqrt{\ysrc^2+4\thetaBH^2})$,
$\theta_+\theta_-=-\thetaBH^2$ and $\theta_++\theta_-=\ysrc$.  So
$|\theta_+|>\thetaBH>|\theta_-|$, and $\mu_+>0$, $\mu_-<0$.
Substituting into $\mu(\theta)$ and writing $u=\ysrc/\thetaBH$ gives
the standard result (eq.~\ref{eq:pm_magnification} of
Chapter~\ref{ch:axisymmetric_models}),
\begin{equation*}
    |\mu_\pm|=\frac{u^2+2}{2u\sqrt{u^2+4}}\pm\frac12 ,
    \qquad
    |\mu_+|+|\mu_-|=\frac{u^2+2}{u\sqrt{u^2+4}} .
\end{equation*}
This exceeds $1$ for every finite $u$, because
$(u^2+2)^2-u^2(u^2+4)=4>0$; it approaches $1$ only as $u\to\infty$.
The signed magnifications satisfy $\mu_++\mu_-=|\mu_+|-|\mu_-|=1$.
For example, $u=5/3$ gives a total magnification of $1.1011$.
\emph{[Attested: \solnlink{15_PowerLaw_PointMass/problems\_15.wl}{problems\_15.wl}, checks 2a--2g.]}

% ----- Exercise: radial radius -----
\subsection*{Exercise~\ref{ex:pl_radial_radius}: Radial Critical Radius}

The radial eigenvalue is
$\lambda_r=1-\dd\alpha_{\rm PL}/\dd\theta=1-(2-\gp)(\thetaE/\theta)^{\gp-1}$.
Setting it to zero,
\begin{equation*}
    \left(\frac{\thetaE}{\theta}\right)^{\gp-1}=\frac{1}{2-\gp}
    \;\Longrightarrow\;
    \frac{\theta_r}{\thetaE}=(2-\gp)^{1/(\gp-1)} .
\end{equation*}
This is real and positive only for $\gp<2$; for $\gp>2$ the base
$2-\gp$ is negative and a non-integer power of it is not real (indeed
$\lambda_r\ge1$ there).  As $\gp\to2^-$ the base tends to $0^+$ with
exponent $\to1$, so $\theta_r\to0$.  At $\gp=1.6$:
$\theta_r/\thetaE=0.4^{1/0.6}=0.4^{5/3}=0.2172$.  The sign of
$\lambda_r$ indeed changes across this radius (negative at
$0.99\,\theta_r$, positive at $1.01\,\theta_r$); likewise at $\gp=1.5$
it changes between $\theta=0.2$ and $0.3$, bracketing $\theta_r=0.25$.
\emph{[Attested: \solnlink{15_PowerLaw_PointMass/problems\_15.wl}{problems\_15.wl}, checks 3a--3e.]}

% ----- Exercise: SIS + PM shifts -----
\subsection*{Exercise~\ref{ex:pl_sis_pm_shift}: SIS plus Point Mass, First-Order Shifts}

Write $\theta_+^{(0)}=\ysrc+\thetaE$ and $\thetaBH^2=\mBH\thetaE^2$.  Expanding the square root,
\begin{equation*}
    \sqrt{(\theta_+^{(0)})^2+4\thetaBH^2}
    =\theta_+^{(0)}\left[1+\frac{2\thetaBH^2}{(\theta_+^{(0)})^2}+\order{\mBH^2}\right] ,
\end{equation*}
so eq.~\eqref{eq:pl_sis_pm_plus} becomes
\begin{equation*}
    \theta_+=\theta_+^{(0)}+\frac{\thetaBH^2}{\theta_+^{(0)}}+\order{\mBH^2} .
\end{equation*}
On the far side, $\theta_-=\tfrac12[(\ysrc-\thetaE)-\sqrt{(\ysrc-\thetaE)^2+4\thetaBH^2}]$;
since $\ysrc-\thetaE<0$, $\sqrt{(\ysrc-\thetaE)^2}=\thetaE-\ysrc$ and the
same expansion gives
\begin{equation*}
    \theta_-=\theta_-^{(0)}+\frac{\thetaBH^2}{\theta_-^{(0)}}+\order{\mBH^2} ,
    \qquad \theta_-^{(0)}=\ysrc-\thetaE<0 .
\end{equation*}
So $\delta\theta_\pm=\thetaBH^2/\theta_\pm^{(0)}$: positive for the
outer image and negative for the counter-image, i.e.\ both move away
from the centre.  For $\thetaE=1$, $\thetaBH=0.08$ ($\mBH=0.0064$),
$\ysrc=0.3$: the exact images are $\theta_+=1.304905$ and
$\theta_-=-0.709026$, against the linearized $1.3+0.0064/1.3=1.304923$
and $-0.7-0.0064/0.7=-0.709143$.  The differences are second order in
$\mBH$.
\emph{[Attested: \solnlink{15_PowerLaw_PointMass/problems\_15.wl}{problems\_15.wl}, checks 4a--4f.]}

% ----- Exercise: tangential shift -----
\subsection*{Exercise~\ref{ex:pl_tangential_shift}: Tangential Critical Curve with a Point Mass}

With $u=\theta/\thetaE$, $\lambda_t(u)=1-u^{1-\gp}-\mBH u^{-2}$.  Put
$u=1+\delta$ and expand to first order in $\delta$ and $\mBH$:
$u^{1-\gp}\simeq1-(\gp-1)\delta$ and $\mBH u^{-2}\simeq\mBH$, so
\begin{equation*}
    \lambda_t\simeq(\gp-1)\,\delta-\mBH=0
    \;\Longrightarrow\;
    \delta=\frac{\mBH}{\gp-1} ,
    \qquad
    \theta_t\simeq\thetaE\left(1+\frac{\mBH}{\gp-1}\right) .
\end{equation*}
The shift grows as the profile becomes shallower ($\gp\to1$), because
then $\bar\kappa$ varies slowly near $\thetaE$ and a small extra mass
moves the radius where $\bar\kappa+\mBH u^{-2}=1$ a long way.  For the
SIS the condition $1-1/u-\mBH/u^2=0$ is a quadratic with root
$u_t=\tfrac12(1+\sqrt{1+4\mBH})=1+\mBH-\mBH^2+\cdots$, in agreement
with the linear result at $\gp=2$.  Numerically, for $\gp=1.5$ and
$\mBH=0.01$, $u_t=1.0195$, against the linear estimate $1.02$.
\emph{[Attested: \solnlink{15_PowerLaw_PointMass/problems\_15.wl}{problems\_15.wl}, checks 5a--5d.]}

% ----- Exercise: annihilation at 3/2 -----
\subsection*{Exercise~\ref{ex:pl_annihilation_15}: Annihilation at $\gp=3/2$ by Hand}

At $\gp=3/2$, $\lambda_r=1-\tfrac12u^{-1/2}+\mBH u^{-2}$.  Multiply by
$u^2$ and put $u=s^2$:
\begin{equation*}
    u^2\lambda_r=s^4-\tfrac12 s^3+\mBH\equiv P(s) .
\end{equation*}
Radial critical curves are positive roots of $P$.  The two roots merge
where $P$ and $P'$ vanish together.  $P'(s)=4s^3-\tfrac32 s^2=0$ gives
$s=3/8$, i.e.\ $u_c=9/64$; then $P(3/8)=0$ gives
\begin{equation*}
    \mBH_{\rm crit}=-\left[\left(\tfrac38\right)^4-\tfrac12\left(\tfrac38\right)^3\right]
    =\frac{27}{4096}\approx0.0066 .
\end{equation*}
Both agree with the general formulae for $u_c$ and
eq.~\eqref{eq:pl_mcrit} at $\gp=3/2$.  The merged radial caustic is at
\begin{equation*}
    \ysrc(u_c)=u_c-\sqrt{u_c}-\frac{\mBH_{\rm crit}}{u_c}
    =\frac{9}{64}-\frac{3}{8}-\frac{3}{64}=-\frac{9}{32} ,
\end{equation*}
i.e.\ $|\ysrc|=9\thetaE/32$.  For comparison, $P$ has two positive
roots at $\mBH=0.003$ and none at $\mBH=0.01$.
\emph{[Attested: \solnlink{15_PowerLaw_PointMass/problems\_15.wl}{problems\_15.wl}, checks 6a--6f.]}

% ----- Exercise: bPL outside -----
\subsection*{Exercise~\ref{ex:bpl_outside}: Broken Power Law outside the Break}

\textbf{(a)} The mass inside the break is
\begin{equation*}
    2\pi\int_0^{\theta_B}\kappa_B\left(\frac{\theta_B}{\theta'}\right)^{t_1}\theta'\,\dd\theta'
    =2\pi\kappa_B\theta_B^{t_1}\,\frac{\theta_B^{2-t_1}}{2-t_1}
    =\frac{2\pi\kappa_B\theta_B^2}{2-t_1} ,
\end{equation*}
and the annulus contributes
$2\pi\kappa_B\theta_B^2[(\theta/\theta_B)^{2-t_2}-1]/(2-t_2)$.  Dividing
the total by $\pi\theta$ gives eq.~\eqref{eq:bpl_alpha_out}.

\textbf{(b)} At $\theta=\theta_B$ the annulus term vanishes and
$\alpha=2\kappa_B\theta_B/(2-t_1)$, which is also the value
of~\eqref{eq:bpl_alpha_in} there.

\textbf{(c)} Split the bracket:
\begin{equation*}
    \alpha=\frac{2\kappa_B\theta_B^{t_2}}{2-t_2}\,\theta^{1-t_2}
      +\frac{2\kappa_B\theta_B^2}{\theta}\left[\frac{1}{2-t_1}-\frac{1}{2-t_2}\right] ,
\end{equation*}
and use $1/(2-t_1)-1/(2-t_2)=(t_1-t_2)/[(2-t_1)(2-t_2)]$.  The first
term is the outer power law continued to the centre; the second is a
point mass with
$\theta_{\rm eff}^2=2\kappa_B\theta_B^2(t_1-t_2)/[(2-t_1)(2-t_2)]$.
\emph{[Attested: \solnlink{15_PowerLaw_PointMass/problems\_15.wl}{problems\_15.wl}, checks B1a--B1e.]}

% ----- Exercise: bPL PM limit -----
\subsection*{Exercise~\ref{ex:bpl_point_mass_limit}: The Point-Mass Limit}

Fixing $2\pi\kappa_B\theta_B^2/(2-t_1)=\pi\thetaBH^2$ means
$\kappa_B=(2-t_1)\thetaBH^2/(2\theta_B^2)$.  Substituting
into~\eqref{eq:bpl_alpha_in},
\begin{equation*}
    \alpha=\frac{2\theta_B^2}{(2-t_1)\theta}\cdot\frac{(2-t_1)\thetaBH^2}{2\theta_B^2}
      \left(\frac{\theta}{\theta_B}\right)^{2-t_1}
    =\frac{\thetaBH^2}{\theta}\left(\frac{\theta}{\theta_B}\right)^{2-t_1}
    \;\xrightarrow{\;t_1\to2\;}\;\frac{\thetaBH^2}{\theta} .
\end{equation*}
The convergence $\kappa_B(\theta_B/\theta)^{t_1}$ becomes more and
more concentrated, but its normalization $\kappa_B\propto(2-t_1)$ goes
to zero: the mass migrates onto the centre while its total stays
fixed.  In the same limit
$\theta_{\rm eff}^2=\thetaBH^2(t_1-t_2)/(2-t_2)\to\thetaBH^2$: the
effective point mass becomes the whole inner mass, as
eq.~\eqref{eq:bpl_eff_ratio} says.
\emph{[Attested: \solnlink{15_PowerLaw_PointMass/problems\_15.wl}{problems\_15.wl}, checks B2a--B2d.]}

% ----- Exercise: bPL SIS limits -----
\subsection*{Exercise~\ref{ex:bpl_sis_limit}: Isothermal Limits}

\textbf{(a)} With $t_1=t_2=1$, eq.~\eqref{eq:bpl_alpha_in} gives
$\alpha=2\kappa_B\theta_B^2\theta^{-1}(\theta/\theta_B)=2\kappa_B\theta_B$,
and in eq.~\eqref{eq:bpl_alpha_out} the bracket is
$1+(\theta/\theta_B-1)=\theta/\theta_B$, giving the same constant.  A
constant deflection is the SIS, with $\thetaE=2\kappa_B\theta_B$.

\textbf{(b)} With $t_2=1$, eq.~\eqref{eq:bpl_identity} reads
\begin{equation*}
    \alpha=2\kappa_B\theta_B+\frac{\theta_{\rm eff}^2}{\theta} ,
    \qquad
    \theta_{\rm eff}^2=\frac{2\kappa_B\theta_B^2(t_1-1)}{2-t_1} ,
    \qquad\theta>\theta_B :
\end{equation*}
an SIS of Einstein radius $\theta_{E,\rm S}=2\kappa_B\theta_B$ plus a
point mass.  Eq.~\eqref{eq:pl_sis_pm_plus} with
$\thetaE\to\theta_{E,\rm S}$ and $\thetaBH^2\to\theta_{\rm eff}^2$
gives the outer image, valid provided it lies outside $\theta_B$.  For
$\theta_B=0.1$, $\kappa_B=5$, $t_1=3/2$: $\theta_{E,\rm S}=1$,
$\theta_{\rm eff}^2=0.1$, and for $\ysrc=0.3$
\begin{equation*}
    \theta_+=\tfrac12\Big[1.3+\sqrt{1.69+0.4}\Big]=1.3728>\theta_B ,
\end{equation*}
which agrees with a direct numerical solution of the broken-power-law
lens equation.

\textbf{(c)} For $t_1<1$ the inner profile is shallower than
isothermal, $\theta_{\rm eff}^2<0$, and the effective point mass is
\emph{negative}: a core removes mass relative to the continued SIS.
Outside the break the lens then looks like a SIS minus a point mass.
\emph{[Attested: \solnlink{15_PowerLaw_PointMass/problems\_15.wl}{problems\_15.wl}, checks B3a--B3d.]}
